\begin{textsample}{NEG Dim 3 – global\_south\_2023\_12 – Score 1.00 – global\_south\_2023\_12/104097194.txt}  \label{ex:f3_neg_010}
From Savarkar to Golwalkar, why Hindutva admires Zionism The foundation of this emerging alliance rests on hostility to Islam, but there are other affinities between Zionism and Hindutva.
Both are characterised by ethno-nationalist ideologies that prioritise factors like race, territory and nativism The crucial consideration was a long-desired national fate in which the state not only aligned with racial identity but also embraced the oversight of racial governance as its fundamental historical mission.
From the Nehruvian national consensus on Palestine to abstaining on a resolution -- the only country in South Asia -- in the United Nations for a humanitarian truce in the"expanded"ground invasion in Gaza, India ' s foreign policy shift is remarkable.
The rapprochement between India and Israel is not new.
Ever since 1992, India ' s strategic and formal diplomatic relations with Israel have increased.
But it was under Prime Minister Modi that, from being a hesitant friend, it became a closer partner of Israel.
Bilateral relations, trade ties, technological assistance, military procurement and counter-terrorism cooperation seem to have acquired political and ideological facets @ @ @ @ @ @ @ @ @ @ in the long history of Hindutva ' s admiration of Zionist ethno-nationalism, possibly because both have found a common enemy in their country ' s largest religious minority.
Hindu nationalists held a fascination with Jews and Zionism from day one as V D Savarkar expressed deep sympathy for a more comprehensive understanding of the Jewish race and underlined its appropriateness for the Hindutva ideology.
He was not only inspired by Zionism and praised illegal colonial settlement but strongly believed that Hindus and Jews shared a history of oppression at the hands of Muslims and both deserved redress.
Underlying his support for the Zionist cause, in Hindutva : Who is a Hindu? ( 1923 ), he wrote :"If the Zionists ' dreams are ever realised -- if Palestine becomes a Jewish state -- it will gladden us almost as much as our Jewish friends."Savarkar cast Jews and Hindus as contenders for claiming their right to national self-determination and the creation of sovereign states by analysing them through the European lens of race.
He asserted that both Hindus and Jews @ @ @ @ @ @ @ @ @ @ groups.
Advertisement Savarkar supported the creation of Israel on both moral and political grounds.
In late 1947, he categorically condemned the vote of the Indian delegation at the UN General Assembly against the plan to divide Palestine into a larger Jewish state and a smaller Arab state.
Instead, he advocated for a binational Arab-Jewish state in Palestine.
The leading RSS ideologue of the time, M S Golwalkar, went further.
In the late 1930s, on the brink of World War II, he wrote,"The Jews had maintained their race, religion, culture and language, and all they wanted was their natural territory to complete their nationality."He displayed a qualified empathy for Zionism and vehemently argued that the only way to atone for past failings would be to gain total authority over a territory and a state.
A significant portion of this intercultural bond resulted from the nationalist and racial insecurity that became the founding principle of the Hindutva ideology.
This, in turn, led to the need @ @ @ @ @ @ @ @ @ @ The foundation of this emerging alliance rests on hostility to Islam, but there are other affinities between Zionism and Hindutva.
Both are characterised by ethno-nationalist ideologies that prioritise factors like race, of course, but also territory and nativism.
While Judaism and Hinduism hold importance for Zionism and Hindutva, these movements are not primarily religious in nature.
For them, Jews and Hindus are not communities of believers, but two peoples.
Interestingly, Savarkar presents Hindus as a"race-jati"in his book.
The other thing Zionism and Hindutva have in common is their view of the territory of their respective countries : Israel is the sacred land of the Jews while India is the punyabhumi of the Hindus.
Ironically, Savarkar also expressed admiration for Adolf Hitler and his policies regarding Jews ( before the Holocaust ).
He commended the anti-Jewish legislation implemented by the Nazis, viewing it as a method to integrate the Jewish minority into the broader Aryan majority in Germany.
This led him to contemplate a parallel situation for Muslims @ @ @ @ @ @ @ @ @ @.
Jews did not matter to him then.
The crucial consideration was a long-desired national fate in which the state not only aligned with racial identity but also embraced the oversight of racial governance as its fundamental historical mission.
Advertisement Jaffrelot is a senior research fellow at CERI-Sciences Po/CNRS, Paris, professor of Indian Politics and Sociology at King ' s India Institute, London.
Joshi is an independent research scholar, journalist-writer and an alumnus of The London School of Economics and Political Science

% matched lemmas: 
\end{textsample}
